% !PRINTED 299
% ?? uncertain: no folio/running head visible on any of the three bands of this page (section-opening page); number inferred from the printed "300" on the immediately following page.
\weberpage{299}

\begin{center}
\gesperrt{Neunter Abschnitt.}

\gesperrt{Abschätzung der Wurzeln.}

\rule{0.3\linewidth}{0.4pt}

\S.\,99.

\gesperrt{Das Budan-Fourier'sche Theorem.}
\end{center}

Die Functionen der \gesperrt{Sturm'schen} Ketten lösen zwar vollständig und ausnahmslos das Problem der Bestimmung der Anzahl der Wurzeln einer Gleichung zwischen gegebenen Grenzen; aber ihre wirkliche Berechnung ist meist schwierig und oft unausführbar. Man kennt eine Reihe von Sätzen zur \gesperrt{Abschätzung} der Zahl der Wurzeln zwischen gegebenen Grenzen, die viel einfachere Regeln liefern, aber freilich auch die Frage nicht vollständig beantworten, sondern nur eine Maximalzahl geben, worüber die Anzahl der Wurzeln nicht hinausgehen kann. Diese Regeln sind oft für die Anwendung sehr nützlich und ausreichend, und sie dürfen daher hier nicht fehlen.

Wir betrachten zunächst ein Verfahren, das von \gesperrt{Budan} herrührt, dann aber von \gesperrt{Fourier} benutzt und erweitert wurde\,\textsuperscript{1}).

Betrachten wir an Stelle einer \gesperrt{Sturm'schen} Kette die Reihe der Ableitungen einer reellen Function $f(x)$ vom $n^{\mathrm{ten}}$ Grade

\weq{1}{f(x),\ f'(x),\ f''(x)\ .\ .\ .\ f^{(n)}(x),}

worin also $f^{(n)}(x)$ eine Constante ist, die wir von Null verschieden und positiv voraussetzen. Sie ist, von einem positiven Zahlenfactor abgesehen, der Co\"efficient von $x^n$ in $f(x)$.

\noindent\rule{0.25\linewidth}{0.4pt}

\footnotesize
\textsuperscript{1}) \gesperrt{Budan}, Nouvelle methode pour la r\'esolution des \'equations num\'eriques, (1803 der Pariser Akademie vorgelegt). \gesperrt{Fourier}, Analyse des \'equations d\'etermin\'ees. Paris 1831. Ueber das Geschichtliche dieser Frage vergleiche man \gesperrt{Lagrange}, Trait\'e de la r\'esolution des \'equations, Note VIII (Werke Bd. 8).
\normalsize
